\documentclass[11pt, reqno]{amsart}

\usepackage[spanish]{babel}
\usepackage[LGR, T1]{fontenc}
\usepackage[utf8]{inputenc}

../LaTeX/general.tex
% \input{../graphics.tex}

\makeatletter
\def\emailaddrname{\textit{Correo electrónico}}
\def\subtitle#1{\gdef\@subtitle{#1}}
\def\@subtitle{}

% Metadata
\def\logo#1{\gdef\@logo{#1}}
\def\@logo{}
\def\institution#1{\gdef\@institution{#1}}
\def\@institution{}
\def\department#1{\gdef\@department{#1}}
\def\@department{}
\def\professor#1{\gdef\@professor{#1}}
\def\@professor{}
\def\course#1{\gdef\@course{#1}}
\def\@course{}
\def\coursecode#1{\gdef\@coursecode{#1}}
\def\@coursecode{}

\renewcommand{\maketitle}{
\begin{center}
	\small
	\renewcommand{\arraystretch}{1.2}
	\begin{tabular}{cp{.37\textwidth}p{0.44\textwidth}}
		% \hline
		\multirow{5}{*}{\includegraphics[height=2.0cm]{\@logo}}
	  & \multicolumn{2}{c}{ \makecell{{\bfseries \@institution} \\ \@department} } \\
	  % & \multicolumn{2}{c|}{{\bfseries\@institution} \\ \@department} \\
	  \cline{2-3}
	  & \textbf{Profesor:} \@professor & \textbf{Ayudante:} \authors \\
	  % \cline{2-3}
	  & \textbf{Curso:} \@course & \textbf{Sigla:} \@coursecode \\
	  % \cline{2-3}
	  & \multicolumn{2}{l}{ \textbf{Fecha:} \@date } \\
	  % \hline
	\end{tabular}
	\\[\baselineskip]
	% {}
	% \vspace{2\baselineskip}
	{\bfseries\Large\@title}
	\ifx\@subtitle\@empty\else
		\\[1ex]
		\large\mdseries\@subtitle
	\fi
\end{center}
}
\makeatother

\usepackage{multirow, makecell}

\usepackage[
	reversemp,
	letterpaper,
	% marginpar=2cm,
	% marginsep=1pt,
	margin=2.3cm
]{geometry}
\usepackage{fontawesome}
% \makeatletter
% \@reversemargintrue
% \makeatother

% Símbolos al margen, necesitan doble compilación
\newcommand{\hard}{\marginnote{\faFire}}
\newcommand{\hhard}{\marginnote{\faFire\faFire}}

% Dependencias para los teoremas
\usepackage{xifthen}
\def\@thmdep{}
\newcommand{\thmdep}[1]{
	\ifthenelse{\isempty{#1}}
	{\def\@thmdep{}}
	{\def\@thmdep{ (#1)}}
}
\newcommand{\thmstyle}{\color{thm}\sffamily\bfseries}

% ===== Estilos de Teoremas ==========
\newtheoremstyle{axiomstyle}
	{0.3cm}
	{0.3cm}
	{\normalfont}
	{0.5cm}
	{\bfseries\scshape}
	{:}
	{4pt}
	{\thmname{#1}\thmnote{ #3}\thmnumber{ (#2)}}
\newtheoremstyle{styleC}
	{0.5cm}
	{0.5cm}
	{\normalfont}
	{0.5cm}
	{\bfseries}
	{:}
	{4pt}
	{\thmname{#1\textrm{\@thmdep}}\thmnumber{ #2}\thmnote{ (#3)}}

% ====== Teoremas (sin borde) ===========
\theoremstyle{axiomstyle}
\newtheorem*{axiom}{Axioma}

% ====== Teoremas (sin borde) ==================
\theoremstyle{styleC}
\newtheorem{thm}{Teorema}[section]
\newtheorem{mydef}[thm]{Definición}
\newtheorem{prop}[thm]{Proposición}
\newtheorem{cor}[thm]{Corolario}
\newtheorem{lem}[thm]{Lema}
\newtheorem{con}[thm]{Conjetura}
\newtheorem*{sol}{Solución}
\newtheorem*{obs}{Observación}
\newtheorem*{ex}{Ejemplo}

\newtcbox{bluebox}[1][]{enhanced jigsaw, 
  sharp corners,
  frame hidden,
  nobeforeafter,
  listing only,
  #1} % comando para crear cajas de colores

\expandafter\let\expandafter\oldproof\csname\string\proof\endcsname
\let\oldendproof\endproof
\renewenvironment{proof}[1][\proofname]{%
  \oldproof[\scshape Demostración:]%
}{\oldendproof} % comando para redefinir la caja de la demostración
\newenvironment{hint}[1][\proofname]{%
  \oldproof[\scshape Pista:]%
}{\oldendproof} % comando para redefinir la caja de la demostración

% colores utilizados
\definecolor{numchap}{RGB}{249,133,29}
\definecolor{chap}{RGB}{6,129,204}
\definecolor{sec}{RGB}{204,0,0}
\definecolor{thm}{RGB}{106,176,240}
\definecolor{thmB}{RGB}{32,31,31}
\definecolor{part}{RGB}{212,66,66}

% ====== Diseño de los titulares ===============
\usepackage[explicit]{titlesec} % para personalizar el documento, la opción <<explicit>> hace que el texto de los titulares sea un objeto interactuable

\titleformat{\subsection}[runin]
	{\bfseries}
	{\textrm{\S}\thesubsection}
	{1ex}
	{#1.}

\setlist[enumerate,1]{label=\arabic*., ref=\arabic*} % Enumerate standards

\usepackage{tikz}
\usetikzlibrary{babel,cd}
% \DeclareMathOperator\Gr{Gr}

\DeclareMathOperator\Cono{Cono}

\title{Homotopías}
\date{\DTMdate{2026-08-27}}

\author{José Cuevas Barrientos}
\email{josecuevasbtos@uc.cl}
\urladdr{https://josecuevas.xyz/teach/2026-2-ayud/}

\logo{../puc_negro.png}
\institution{Pontificia Universidad Católica de Chile}
\department{Facultad de Matemáticas}
\course{Topología Algebraica}
\coursecode{MAT2850}
\professor{Mauricio Bustamante}

\begin{document}

\maketitle

\section{Homotopía}
\begin{mydef}
	Sean $\varphi_\bullet, \psi_\bullet \colon A_\bullet \to B_\bullet$ un par de homomorfismos de complejos.
	Una \strong{homotopía} es una familia de homomorfismos $s_n \colon A_n \to B_{n+1}$ tales que
	\[
		\varphi_n - \psi_n = \ud^B_{n+1}\circ s_n + s_{n-1}\circ \ud^A_n,
	\]
	en cuyo caso, denotamos que $s \colon \varphi_\bullet \simeq \psi_\bullet$.

	Se dice que dos complejos $A_\bullet, B_\bullet$ son \emph{homotópicos} si existen homomorfismos $\alpha_\bullet \colon A_\bullet
	\to B_\bullet$ y $\beta_\bullet \colon B_\bullet \to A_\bullet$ tales que $\beta\circ\alpha \simeq \Id_A$ y $\alpha\circ\beta \simeq
	\Id_B$.
\end{mydef}

\begin{enumerate}
	\item Sean $\varphi_\bullet, \psi_\bullet \colon A_\bullet \to B_\bullet$ un par de homomorfismos de complejos que son homólogos.
		Pruebe que $H_n(\varphi) = H_n(\psi) \colon H_n(A) \to H_n(B)$ para todo $n \in \Z$.

	\item Pruebe que un complejo $A_\bullet$ es homótopico al nulo $0_\bullet$ syss es exacto y escindido, es decir, para todo $n$ la
		sucesión exacta $0 \to Z_n(A) \to A_n \to \ud A_n \to 0$ se escinde.

	\item Sea $\varphi_\bullet \colon A_\bullet \to B_\bullet$ un homomorfismo de complejos y defina, para cada $n \in \Z$ el módulo $\Cono_n(\varphi) := A_{n-1} \oplus B_n$ con los diferenciales
		\[
			\ud =
			\begin{bsmallmatrix}
				-\ud^A & 0 \\
				-\varphi_{n-1} & \ud^B
			\end{bsmallmatrix} 
			\colon \Cono_n(\varphi) \longrightarrow \Cono_{n-1}(\varphi), \qquad (a, b) \longmapsto (-\ud a, \ud b - \varphi a).
		\]
		\begin{enumerate}
			\item Pruebe que $\Cono_\bullet(\varphi)$ es un complejo.
			\item Pruebe que 
				\begin{tikzcd}[cramped, sep=small]
					0 \rar & B_\bullet \rar["i^2"] & \Cono_\bullet(\varphi) \rar["\pi^1"] & A_{\bullet-1} \rar & 0
				\end{tikzcd}
				es una sucesión exacta de complejos.
			\item Pruebe que $\varphi_\bullet$ es un cuasi-isomorfismo (i.e., induce un isomorfismo $H_n(\varphi)$ en homología para cada $n$) syss $\Cono_\bullet(\varphi)$ es exacto.
		\end{enumerate}

	\item Sean $F_\bullet \to M$ y $G_\bullet \to N$ un par de resoluciones libres de los $R$-módulos $M$ y $N$.
		\begin{enumerate}
			\item Pruebe que todo homomorfismo $f \colon M \to N$ admite un \emph{levantamiento} $\varphi_\bullet \colon F_\bullet \to G_\bullet$.
				Es decir, un homomorfismo de complejos tal que el siguiente diagrama conmuta
				% https://q.uiver.app/#q=WzAsOCxbMywwLCJNIl0sWzMsMSwiTiJdLFsyLDAsIkZfMCJdLFsyLDEsIkdfMCJdLFsxLDAsIkZfMSJdLFsxLDEsIkdfMSJdLFswLDAsIlxcY2RvdHMiXSxbMCwxLCJcXGNkb3RzIl0sWzAsMSwiZiJdLFszLDEsIiIsMix7InN0eWxlIjp7ImhlYWQiOnsibmFtZSI6ImVwaSJ9fX1dLFsyLDAsIiIsMCx7InN0eWxlIjp7ImhlYWQiOnsibmFtZSI6ImVwaSJ9fX1dLFs2LDRdLFs0LDJdLFs3LDVdLFs1LDNdLFsyLDMsIlxcdmFycGhpXzAiLDAseyJzdHlsZSI6eyJib2R5Ijp7Im5hbWUiOiJkYXNoZWQifX19XSxbNCw1LCJcXHZhcnBoaV8xIiwwLHsic3R5bGUiOnsiYm9keSI6eyJuYW1lIjoiZGFzaGVkIn19fV1d
				\[\begin{tikzcd}
					\cdots & {F_1} & {F_0} & M \\
					\cdots & {G_1} & {G_0} & N
					\arrow[from=1-1, to=1-2]
					\arrow[from=1-2, to=1-3]
					\arrow["{\varphi_1}", dashed, from=1-2, to=2-2]
					\arrow[two heads, from=1-3, to=1-4]
					\arrow["{\varphi_0}", dashed, from=1-3, to=2-3]
					\arrow["f", from=1-4, to=2-4]
					\arrow[from=2-1, to=2-2]
					\arrow[from=2-2, to=2-3]
					\arrow[two heads, from=2-3, to=2-4]
				\end{tikzcd}\]

			\item Pruebe que si $\varphi_\bullet, \psi_\bullet \colon F_\bullet \to G_\bullet$ son levantamientos de un mismo homomorfismo $f \colon M \to N$,
				entonces $\varphi_\bullet \simeq \psi_\bullet$.

			\item Concluya que los funtores $\Ext_R^n(M, N)$ están bien definidos.
		\end{enumerate}
\end{enumerate}

\printbibliography

\end{document}
